\scanpage{109}
\textcircled{2} Cas où on dispose seulement d'une représentation de $F$ sous forme sympathique \uline{localement} sur $X$ pour la top. (Zar). Dans ce cas, \deleted{on ne peut plus utiliser \unsure{comme} l'hyp. $X$ séparé, on doit en tous cas localiser sur $X$ au sens (ét), et utiliser la suite spectrale de Leray d'un recouvrement. On trouve \unsure{au moins} l'effectivité, poids entre $0$ et $2n+2j$}
on a intérêt à travailler \uline{d'abord} avec la cohomologie à support propre sur $X$, qui permet de découper $X$ en morceaux, et d'obtenir le $=$ résultat que dessus. On déduit l'analogue en cohomologie à supports quelconques par dualité.

\deleted{Si $F$ est sympathique,}
\uline{Remarque} [ En fait, on voit par ce raisonnement que, en ce qui concerne la coh. à support propre, les résultats sont obtenus en supposant seulement que $F$ soit \textquotedblleft joli par morceaux\textquotedblright{} sur $X$ ; \inserted{\unsure{la lissité|se lissité} de $X$ \unsure{pas nécessaire|par récurrence|par ailleurs}} ]. Le fait $F$ constant tordu \inserted{sur $X$ lisse} est alors utilisé \deleted{\unsure{dans le cas li}} pour pouvoir procéder par dualité.

\textcircled{C} \uline{Cas de coefficients effectifs de (poids} $j$, \uline{niveau} $j$) \uline{quelc.}, $X$ quelconque.

1) Cas de la cohomologie $H^i_!$ : déjà traité dans la remarque précé\deleted{\unsure{dente}}dente : résultat le $=$ que dans \textcircled{B}.1).

2) Cas de la cohomologie $H^i(X,F)$. Je dis qu'il est \uline{effectif}, de poids \deleted{entre $0$ et $2i$}, \inserted{$\le\operatorname{Inf}(2i,2n)=2\operatorname{Inf}(i,n)$}, de niveau $\le\operatorname{Inf}(i,n)$, où $n=\dim X$. \deleted{On procède par} \ldots
\margin{?}
% archive sheet number: 108
